\documentclass[11pt]{article}
\usepackage[T1]{fontenc}
\usepackage{lmodern}
\usepackage[margin=1in]{geometry}
\usepackage{amsmath,amssymb,amsthm,mathtools}
\usepackage{microtype}
\usepackage{booktabs,tabularx}
\usepackage{xcolor}
\usepackage[colorlinks=true,linkcolor=blue!50!black,urlcolor=blue!50!black,citecolor=blue!50!black]{hyperref}
\newtheorem{theorem}{Theorem}
\newtheorem{lemma}[theorem]{Lemma}
\newcommand{\FL}{\mathrm{FL}}
\newcommand{\FLpos}{\mathrm{FL}_{\geq0}}
\newcommand{\sharpL}{\#\mathrm L}
\newcommand{\TC}{\mathrm{TC}}
\newcommand{\len}{\ell}
\title{Principal Schur specialization and\protect\\ GL2 plethysm in logarithmic space}
\author{\href{https://github.com/fengqiaoqu}{fengqiaoqu}}
\date{22 October 2026}
\hypersetup{pdftitle={Principal Schur specialization and GL2 plethysm in logarithmic space},pdfauthor={fengqiaoqu}}
\begin{document}
\maketitle
\begin{abstract}
For an arbitrary partition $\lambda$ and unary integers $k,r\geq0$, we show
that $[q^r]s_\lambda(1,q,\ldots,q^k)$ is computable in deterministic
logarithmic space, with binary output. The same holds for two-row
$\mathrm{GL}_2$ plethysm coefficients with an arbitrary outer partition,
and for Hermite coefficients. In particular these functions belong to
$\FLpos\subseteq\sharpL$. The proof combines carry-free integer substitution,
the hook-content formula, and the Hesse--Allender--Barrington integer-arithmetic
theorem. It also permits binary encoding of $k$ when the partition and $r$
remain unary. This gives affirmative answers to Questions 6.4, 6.22, and 6.29
of the fixed v1 PDF of \emph{Counting in logarithmic space}. The result is an
application of classical tools; external review and priority assessment
remain outstanding.
\end{abstract}

\section{Statements and conventions}
For a partition $\lambda=(\lambda_1\geq\cdots\geq\lambda_\ell>0)$ of
size $w=|\lambda|$, put
\[
S_{\lambda,k}(q)=s_\lambda(1,q,\ldots,q^k)
=\sum_{t\geq0}c_t(\lambda,k)q^t.
\]
The empty partition is allowed. Initially, the parts and the nonnegative
integers $k,r$ are unary, with fixed separators; $L$ denotes total input
length. Invalid encodings have value zero. There is no constant bound on
$\len(\lambda)$. Outputs are binary, without leading zeros except for the
output $0$.

We use $\FL$ for deterministic transducers with read-only input,
$O(\log L)$ work space, and write-only output. All transducers constructed
below have polynomial time and output length. The notation $\FLpos$
restricts to nonnegative integer-valued functions. The class $\sharpL$
counts accepting paths of polynomial-time nondeterministic logarithmic-space
machines.

\begin{theorem}\label{thm:main}
The unary-input function $(\lambda,k,r)\mapsto c_r(\lambda,k)$ belongs to
$\FLpos\subseteq\sharpL$.
\end{theorem}

For $\mu\vdash w$, define plethysm coefficients by
$s_\mu[s_{(k)}]=\sum_{\nu\vdash wk}a_{\mu[(k)]}^{\nu}s_\nu$.
\begin{theorem}\label{thm:pleth}
For unary input and $0\leq2r\leq wk$, the function
$(\mu,k,r)\mapsto a_{\mu[(k)]}^{(wk-r,r)}$ belongs to
$\FLpos\subseteq\sharpL$, without a bound on $\len(\mu)$.
Assign value zero outside this range.
\end{theorem}

\begin{theorem}\label{thm:mixed}
The conclusions of Theorems~\ref{thm:main} and \ref{thm:pleth} remain true
if $k$ is binary and the partition and $r$ remain unary. In
Theorem~\ref{thm:pleth}, the target shape is specified implicitly by
$(\mu,k,r)$.
\end{theorem}

\section{Arithmetic and virtual inputs}
Our arithmetic input is the theorem of Hesse, Allender, and
Barrington~\cite[Corollary 6.3]{HAB}: integer division and iterated
multiplication are in DLOGTIME-uniform $\TC^0$. In particular, a
polynomial number of explicitly represented binary factors of polynomial
total length can be multiplied, and two such binary integers can be
divided, by deterministic logarithmic-space bit algorithms.

To see the space consequence, evaluate a single output of a uniform
constant-depth threshold circuit by recursive recomputation. Each gate
level needs a gate index, an input index, and a counter of satisfied
inputs; the latter computes a threshold gate. Uniformity generates the
wiring and gate types. There are only constantly many levels and each
index or counter uses logarithmic space in the arithmetic input length.
The circuit has polynomial size and constant depth, so recomputation
also takes polynomial time. Repeating for each output bit gives a
transducer.

The 2014 corrigendum~\cite{HABerr} retracts parts 1 and 3 of Corollary 6.6,
which concern bounded arithmetic. It does not retract the
integer-arithmetic theorem used here. Neither those retracted statements
nor the polynomial-arithmetic corollary is needed.

\begin{lemma}[Virtual-input composition]\label{lem:compose}
A constant number of logarithmic-space transducers with polynomially
bounded intermediate lengths can be composed in logarithmic space and
polynomial time.
\end{lemma}
\begin{proof}
To simulate $F(G(x))$, store the virtual input-head position $j$ of $F$.
When $F$ reads position $j$, rerun $G$ using an output-position counter
until the requested bit is reached. Polynomially bounded positions need
$O(\log L)$ bits. The simulation state and the generating subroutine
use additive work space. Their times and the number of recomputations
are polynomial. A fixed number of such layers preserves these bounds,
without storing an intermediate output in full.
\end{proof}

\begin{lemma}[Carry-free substitution]\label{lem:blocks}
Suppose $P(q)=\sum_{t=0}^d u_tq^t$ and $0\leq u_t<2^W$. For $B=2^W$,
\begin{equation}\label{eq:block}
u_r=\left\lfloor\frac{P(B)}{2^{Wr}}\right\rfloor\bmod2^W.
\end{equation}
If the relevant lengths and positions are polynomially bounded and the
bit function of $P(B)$ is computable in logarithmic space, so is the
binary output $u_r$.
\end{lemma}
\begin{proof}
The sum $\sum_tu_tB^t$ is already a base-$B$ expansion. Thus $u_r$ is
the width-$W$ binary block starting at position $Wr$, numbering the
least significant bit by zero. Scan this block to find its highest
nonzero bit, then output downwards from that position, or output zero
if no such bit exists. Only counters and the bit algorithm are needed.
\end{proof}

\section{Principal specialization}
\begin{proof}[Proof of Theorem~\ref{thm:main}]
Repeated scans with counters validate the input and compute $w,k,r$
and $\len(\lambda)$. If $w=0$, then $S_{\lambda,k}=1$. If
$\len(\lambda)>k+1$, then $S_{\lambda,k}=0$. If $r>wk$, its coefficient
is zero. These boundary cases are immediate. Otherwise $w>0$ and,
with $K=k+1$, we have $\len(\lambda)\leq K$. Set
\begin{equation}\label{eq:params}
W=1+w\lceil\log_2K\rceil,\qquad B=2^W,\qquad
b(\lambda)=\sum_i(i-1)\lambda_i.
\end{equation}

\paragraph{A strict coefficient bound.}
The tableau expansion of a Schur polynomial has nonnegative
coefficients. Its value at $q=1$ counts semistandard tableaux with
$K$ available entries. Counting all assignments to the $w$ boxes gives
\begin{equation}\label{eq:bound}
\sum_t c_t(\lambda,k)=s_\lambda(1^K)
\leq K^w\leq2^{w\lceil\log_2K\rceil}<B.
\end{equation}
This does not assume a length bound. When $K=1$, it gives $W=1$ and
still holds. Every coefficient therefore satisfies
Lemma~\ref{lem:blocks}.

\paragraph{Exact integer evaluation.}
For each box $u=(i,j)\in[\lambda]$, put
\[
a_u=K+j-i,\qquad h_u=\lambda_i-j+\lambda'_j-i+1.
\]
The principal-specialization hook-content formula, also used in
\cite[\S6]{GIP}, states
\[
S_{\lambda,k}(q)=q^{b(\lambda)}
\frac{\prod_u(1-q^{a_u})}{\prod_u(1-q^{h_u})}.
\]
Since $\len(\lambda)\leq K$, the exponents satisfy
\begin{equation}\label{eq:positive}
1\leq a_u\leq K+w,\qquad1\leq h_u\leq w.
\end{equation}
At $q=B>1$, cancel the signs of the $w$ numerator and $w$ denominator
factors to obtain
\begin{equation}\label{eq:quotient}
S_{\lambda,k}(B)=
\frac{B^{b(\lambda)}\prod_u(B^{a_u}-1)}
{\prod_u(B^{h_u}-1)}.
\end{equation}
The denominator is positive, and the quotient is an exact integer:
it is the value of an integer-coefficient polynomial at an integer.
Only integer multiplication and division are required.

\paragraph{Polynomial binary length.}
Unary encoding gives $w,k,r=O(L)$. Also $b(\lambda)\leq w^2$ and,
using a deliberately coarse bound, $W=O((L+1)^2)$. The numerator in
\eqref{eq:quotient} has bit length at most
\begin{align}\label{eq:length}
1+W\left(b(\lambda)+\sum_u a_u\right)
&\leq1+W\bigl(w^2+w(K+w)\bigr)\notag\\
&=O((L+1)^4).
\end{align}
The denominator has at most $1+Ww^2$ bits; the quotient is no longer
than the numerator. The integers may be exponentially large in value,
but their binary representations have polynomial length.

\paragraph{Generating factors and extracting a block.}
A stored box coordinate $(i,j)$ and a scan determine $\lambda_i$.
Another scan counts the parts at least $j$ to obtain $\lambda'_j$.
All exponents, positions, and quantities in \eqref{eq:params} and
\eqref{eq:positive} are polynomially bounded and computable with a
constant number of logarithmic-size counters.

Each factor $B^e-1=2^{We}-1$ is a string of $We$ ones.
The factor $B^{b(\lambda)}$ is a one followed by $Wb(\lambda)$ zeros.
Consequently, the factor lists in \eqref{eq:quotient}, including padding
and separators for a fixed arithmetic encoding, are polynomial-length
virtual inputs with logarithmic-space bit functions. Apply iterated
multiplication separately to the two lists and integer division to the
products. By Lemma~\ref{lem:compose}, the bit function of
$S_{\lambda,k}(B)$ is in logarithmic space. There are constantly many
composition layers: the number of boxes is the number of factors in
one iterated-multiplication instance, not the number of layers.
Finally \eqref{eq:bound} and Lemma~\ref{lem:blocks} compute $c_r$ in $\FL$.

\paragraph{A counting machine.}
For the $\sharpL$ inclusion, use $0\leq c_r<2^W$. Guess exactly $W$
bits, most significant first, specifying $y\in\{0,\ldots,2^W-1\}$.
Maintain a position counter and the three-way comparison state of the
prefixes of $y$ and $c_r$. At each position, deterministically recompute
the corresponding bit of $c_r$, padded with zeros. Accept exactly when
$y<c_r$. There is one path per length-$W$ binary string and hence exactly
$c_r$ accepting paths. The machine has logarithmic work space and
polynomial time on every path. Each guessed bit is consumed once;
past nondeterministic choices are not guessed again to reconstruct a
witness. Boundary cases use the same construction or machines with
zero or one accepting path.
\end{proof}

\section{Plethysm and Hermite coefficients}
\begin{proof}[Proof of Theorem~\ref{thm:pleth}]
Evaluating the plethysm alphabet gives
\[
s_\mu[s_{(k)}](1,q)=s_\mu(1,q,\ldots,q^k),
\]
because $s_{(k)}(1,q)=1+q+\cdots+q^k$. Schur polynomials with more
than two rows vanish in these two variables, while
\[
s_{(wk-t,t)}(1,q)=q^t+q^{t+1}+\cdots+q^{wk-t}
\quad(0\leq2t\leq wk).
\]
For $0\leq2r\leq wk$, comparison of adjacent coefficients gives
\begin{equation}\label{eq:diff}
a_{\mu[(k)]}^{(wk-r,r)}=c_r(\mu,k)-c_{r-1}(\mu,k),\qquad c_{-1}=0.
\end{equation}
Both terms are in $\FL$. Subtraction is in logarithmic space: to output
a bit, rescan less significant positions to recompute its borrow,
keeping a single borrow bit. Compose this with the coefficient bit
algorithms using Lemma~\ref{lem:compose}. The difference is nonnegative
because it is a representation multiplicity, and it is at most
$c_r<2^W$. The fixed-length comparison machine in the preceding proof
puts it in $\sharpL$. No closure of $\sharpL$ under subtraction is used,
and no bound on $\len(\mu)$ was imposed.
\end{proof}

For $\mu=(n)$, let $p_r(n\times k)$ count partitions of $r$ contained
in an $n$-by-$k$ rectangle. Then
\[
s_{(n)}(1,q,\ldots,q^k)=\genfrac{[}{]}{0pt}{}{n+k}{k}_q
=\sum_rp_r(n\times k)q^r.
\]
The Hermite identity \cite[(6.1)--(6.2)]{GIP} yields
\begin{equation}\label{eq:hermite}
a_{(n)[(k)]}^{(nk-r,r)}
=g\bigl((nk-r,r),(n^k),(n^k)\bigr)
=p_r(n\times k)-p_{r-1}(n\times k)
\end{equation}
for $0\leq2r\leq nk$, where $(n^k)$ has $k$ parts equal to $n$.
Empty cases are interpreted by the constant Schur polynomial. These
coefficients are therefore in $\FLpos\subseteq\sharpL$.

There is also a shorter special-case size bound. Set $W=n+k+1$,
as $\sum_rp_r(n\times k)=\binom{n+k}{k}<2^W$. At $q=2^W$,
\[
\genfrac{[}{]}{0pt}{}{n+k}{k}_{2^W}
=\frac{\prod_{i=1}^k(2^{W(n+i)}-1)}
{\prod_{i=1}^k(2^{Wi}-1)}.
\]
For $k>0$, the numerator and denominator have at most
$Wk(2n+k+1)/2$ and $Wk(k+1)/2$ bits, respectively. If $k=0$,
both products are empty and the quotient is one.

\section{A mixed-encoding extension}
\begin{proof}[Proof of Theorem~\ref{thm:mixed}]
For every $0\leq t\leq r$, the tableau expansion shows
\begin{equation}\label{eq:truncate}
[q^t]s_\lambda(1,q,\ldots,q^k)
=[q^t]s_\lambda(1,q,\ldots,q^{\min(k,r)}).
\end{equation}
A tableau using a letter of weight $q^j$ for $j>r$ has total degree
greater than $r$, since all other weights have nonnegative degree.
Tableaux using only retained letters satisfy exactly the same row and
column conditions. This proves the claimed bijective restriction,
including cases where the truncated alphabet is too short for the shape.

While reading binary $k$, maintain a saturated counter updated by
$z\leftarrow\min(r,2z+b)$ for the next bit $b$. This computes
$\min(k,r)$ in logarithmic space. Its virtual unary representation has
polynomial length because $r$ is unary. Theorem~\ref{thm:main} and
Lemma~\ref{lem:compose} apply, proving the coefficient claim.
The same truncation preserves both coefficients in \eqref{eq:diff}.

For the validity condition of Theorem~\ref{thm:pleth}, if $w>0$ compute
$\widehat k=\min(k,2r+1)$ by saturation. Then
$2r\leq wk$ if and only if $2r\leq w\widehat k$, and the latter
comparison uses polynomially bounded integers. If $w=0$, validity
requires $r=0$. Thus the implicit target shape can be checked without
storing a large binary product. This extension does not encode all
parameters in binary.
\end{proof}

\section{Relation to the stated questions and limitations}
The reference for numbering is the fixed arXiv v1 PDF of
Guti\'errez, Ikenmeyer, and Panova~\cite{GIP}, dated 17 July 2026.
Its experimental HTML rendering has different labels:
\begin{center}
\small
\begin{tabularx}{\linewidth}{@{}lllX@{}}
\toprule
v1 PDF question & Page & HTML label & Consequence here\\
\midrule
6.4 & 28 & 6.4 & Hermite coefficients are in $\FLpos\subseteq\sharpL$.\\
6.22 & 36 & 6.21 & Principal specialization, with arbitrary partition length.\\
6.29 & 38 & 6.27 & Two-row $\mathrm{GL}_2$ plethysm, with arbitrary outer length.\\
\bottomrule
\end{tabularx}
\end{center}
For Question 6.29, $\sharpL\subseteq
\#\mathrm{TISP}(\operatorname{poly}(L),\log^2L)$ gives the requested
inclusion. These are three consequences of one main coefficient theorem.
The mixed-encoding extension is not counted as a separately sourced
open question.

The argument uses a single-variable polynomial whose exact integer
evaluation has polynomial binary length. It does not give an algorithm
of this kind for general Littlewood--Richardson or Kronecker coefficients,
and does not resolve whether general Kronecker coefficients belong to
$\#\mathrm P$ or separate them from that class.

The accompanying Python checks compare exact integer substitution with
independent tableau and rectangle-partition computations. They check
finite values and identities; they do not implement the logarithmic-space
multiplication and division transducers. The space theorem follows from
the proof and \cite{HAB}, not from the memory use of those programs.

\paragraph{Research-status notice.}
This draft received internal AI-assisted mathematical review and finite
exact checks. It has not received external peer review or author
confirmation, and is not a Lean formalization. No global priority claim
is made. The central observation applies classical formulas and
integer-arithmetic results to the stated coefficient functions.

\begin{thebibliography}{9}
\bibitem{GIP}
\'Alvaro Guti\'errez, Christian Ikenmeyer, and Greta Panova.
\emph{Counting in logarithmic space}.
arXiv:2607.15881v1, 17 July 2026.
\href{https://arxiv.org/pdf/2607.15881v1}{Fixed v1 PDF};
\href{https://arxiv.org/html/2607.15881v1}{v1 HTML}.

\bibitem{HAB}
William Hesse, Eric Allender, and David A. Mix Barrington.
Uniform constant-depth threshold circuits for division and iterated
multiplication.
\emph{Journal of Computer and System Sciences} \textbf{65}(4),
695--716, 2002.
\href{https://people.cs.rutgers.edu/~allender/papers/division.pdf}{Author PDF}.

\bibitem{HABerr}
William Hesse, Eric Allender, and David A. Mix Barrington.
Corrigendum for Uniform Constant-Depth Threshold Circuits for Division
and Iterated Multiplication.
\emph{Journal of Computer and System Sciences} \textbf{80},
496--497, 2014.
\href{https://people.cs.rutgers.edu/~allender/papers/hab.corrigendum.pdf}{Author PDF}.
\end{thebibliography}
\end{document}
